\RequirePackage{fix-cm}
\documentclass[10pt,twocolumn]{article}

\usepackage[margin=0.6in,top=0.55in,bottom=0.6in]{geometry}
\setlength{\columnsep}{0.22in}
\usepackage{booktabs}
\usepackage{array}
\usepackage{amsmath}
\usepackage{amssymb}
\usepackage{hyperref}
\usepackage{caption}
\usepackage{titlesec}

\titlespacing*{\section}{0pt}{0.5ex}{0.2ex}
\titlespacing*{\subsection}{0pt}{0.4ex}{0.15ex}
\titlespacing*{\paragraph}{0pt}{0.3ex}{0.5em}
\setlength{\parskip}{0.35ex}
\setlength{\parindent}{1em}
\captionsetup{skip=2pt,font=small}
\renewcommand{\arraystretch}{0.95}

\begin{document}

\twocolumn[{%
\begin{center}
{\Large\bfseries Goodreads Recommendation System}\par
\vspace{2pt}
{\small Repository: \url{https://github.com/yoonseol-jang/goodreads-recommendation-system}}
\end{center}
\vspace{1ex}
}]

% =====================================================================
% Shared intro
% =====================================================================
\section*{Preprocessing and Data Splitting}

\paragraph{CSV-to-Parquet Conversion.} The raw interaction file contains approximately 223 million rows. Reading it as CSV on every experiment incurs a full schema-inference pass in addition to the data scan, which is wasteful on a shared cluster. As the first step of the pipeline, we converted the file to Parquet using an explicit schema definition rather than \texttt{inferSchema=True}, avoiding the extra scan entirely. All downstream scripts read from the Parquet version, which benefits from columnar storage, predicate pushdown, and compressed I/O.

\paragraph{User-Based Sampling.} Because full-dataset ALS and MinHash runs were too costly on the shared cluster, the final experiments use a deterministic 3\% user sample. The sample was drawn by computing \texttt{xxhash64(user\_id) \% 100} and keeping users whose hash value falls below 3. We also used the same deterministic hash rule with a 1\% threshold for early hyperparameter and weight screens, then reran selected configurations on the 3\% sample for final comparisons. Because \texttt{xxhash64} is a pure function of \texttt{user\_id}, every row belonging to a kept user passes the filter in a single scan with no additional distinct-and-join shuffle. The resulting sample preserves each sampled user's complete interaction history, maintaining the natural distribution of interaction counts rather than distorting it by downsampling individual records.

\paragraph{Train/Validation/Test Split.} Recommendation evaluation requires that each user's history be divided into an observed portion, used to generate recommendations, and a held-out portion, used to score them. We first discarded users with fewer than five interactions, as they do not provide enough history to support a meaningful three-way split. For each remaining user, interactions were randomly ordered with a fixed seed (\texttt{seed=42}) and then partitioned within each user: the first 60\% were assigned to training, the next 20\% to validation, and the final 20\% to test. Splitting within users rather than across users guarantees that every validation and test user also has a training history, eliminating cold-start cases from the evaluation cohort by construction. All three splits were written to Parquet for reuse across all subsequent experiments.

% =====================================================================
% Section 1: Market Segmentation
% =====================================================================
\section{Market Segmentation}

For market segmentation, we represented each book by the set of users who read it, using only interactions where \texttt{is\_read = 1}. To reduce unstable similarities from very rare books, we filtered out books with fewer than 30 readers in the 3\% user sample before computing similarity. Each book's reader set was then encoded as a binary sparse vector via \texttt{CountVectorizer} (\texttt{binary=True}), and Spark \texttt{MinHashLSH} was fit on these vectors with \texttt{numHashTables=5} to produce signatures that approximate Jaccard similarity in a fixed-length, low-dimensional space.

\paragraph{Scalability.} A brute-force search over all book pairs is $O(N^2)$ in the number of books, which is prohibitive even on the 3\% sample. \texttt{MinHashLSH} avoids this by hashing each signature into buckets such that only books that collide in at least one of the 5 hash tables become candidate pairs --- highly similar books collide with high probability, while dissimilar ones rarely do. We invoked \texttt{approxSimilarityJoin} with a Jaccard distance threshold of $0.8$ (i.e.\ candidates with similarity $\geq 0.2$), which is loose enough to retain all genuinely high-overlap pairs while still discarding the vast majority of the $O(N^2)$ search space at the hashing stage. The shortlist of candidate pairs is small enough that we then compute exact intersection and union sizes per pair in a single distributed pass, so MinHash supplies the recall and the exact step supplies the final ranking precision.

The final top 100 pairs were sorted by exact Jaccard similarity. The highest-ranked pair had 34 shared readers out of 35 users in the union, giving a similarity of 0.9714. Overall, the top 100 pairs had high reader overlap, with similarities ranging from about 0.97 to 0.86, suggesting that the method identified books occupying similar positions in the reading market.

\begin{table*}[t]
\centering
\fontsize{7.5}{8.0}\selectfont
\setlength{\tabcolsep}{1.5pt}
\caption{Top 100 most similar book pairs from the 3\% user sample, after filtering to books with at least 30 readers. Pairs are sorted by exact Jaccard similarity; the left block lists ranks 1--50 and the right block lists ranks 51--100. $r_1, r_2$ are the reader counts of each book; $|\cap|, |\cup|$ are intersection and union sizes of the reader sets.}
\label{tab:minhash}
\begin{tabular}{rrrrrrc@{\hspace{1em}}rrrrrrc}
\toprule
$b_1$ & $b_2$ & $r_1$ & $r_2$ & $|\cap|$ & $|\cup|$ & $J$ & $b_1$ & $b_2$ & $r_1$ & $r_2$ & $|\cap|$ & $|\cup|$ & $J$ \\
\midrule
3952 & 3953 & 35 & 34 & 34 & 35 & 0.971429 & 21846 & 21848 & 50 & 54 & 49 & 55 & 0.890909 \\
21738 & 21739 & 71 & 70 & 69 & 72 & 0.958333 & 17471 & 17473 & 52 & 50 & 48 & 54 & 0.888889 \\
21736 & 21737 & 66 & 67 & 65 & 68 & 0.955882 & 21840 & 21843 & 33 & 35 & 32 & 36 & 0.888889 \\
3958 & 3959 & 44 & 42 & 42 & 44 & 0.954545 & 3951 & 3952 & 33 & 35 & 32 & 36 & 0.888889 \\
4383 & 4385 & 42 & 44 & 42 & 44 & 0.954545 & 21840 & 21841 & 33 & 35 & 32 & 36 & 0.888889 \\
4383 & 4384 & 42 & 42 & 41 & 43 & 0.953488 & 58991 & 58992 & 119 & 117 & 111 & 125 & 0.888 \\
58908 & 58910 & 41 & 43 & 41 & 43 & 0.953488 & 14229 & 83430 & 59 & 58 & 55 & 62 & 0.887097 \\
3955 & 3957 & 39 & 39 & 38 & 40 & 0.95 & 8264 & 8268 & 96 & 102 & 93 & 105 & 0.885714 \\
112530 & 112531 & 34 & 36 & 34 & 36 & 0.944444 & 58921 & 58922 & 72 & 73 & 68 & 77 & 0.883117 \\
207578 & 207579 & 32 & 32 & 31 & 33 & 0.939394 & 13976 & 13977 & 274 & 289 & 264 & 299 & 0.882943 \\
51199 & 67248 & 32 & 32 & 31 & 33 & 0.939394 & 56477 & 56478 & 32 & 32 & 30 & 34 & 0.882353 \\
44736 & 44737 & 31 & 33 & 31 & 33 & 0.939394 & 58934 & 58935 & 31 & 33 & 30 & 34 & 0.882353 \\
12433 & 12434 & 31 & 31 & 30 & 32 & 0.9375 & 140338 & 140339 & 31 & 33 & 30 & 34 & 0.882353 \\
44669 & 44673 & 40 & 43 & 40 & 43 & 0.930233 & 113027 & 113028 & 39 & 40 & 37 & 42 & 0.880952 \\
21737 & 21739 & 67 & 70 & 66 & 71 & 0.929577 & 157997 & 157998 & 31 & 31 & 29 & 33 & 0.878788 \\
58916 & 58918 & 52 & 54 & 51 & 55 & 0.927273 & 5429 & 5430 & 31 & 31 & 29 & 33 & 0.878788 \\
19320 & 19321 & 91 & 90 & 87 & 94 & 0.925532 & 3910 & 3911 & 152 & 163 & 147 & 168 & 0.875 \\
89049 & 89051 & 47 & 47 & 45 & 49 & 0.918367 & 21847 & 21848 & 51 & 54 & 49 & 56 & 0.875 \\
21737 & 21738 & 67 & 71 & 66 & 72 & 0.916667 & 5831 & 5833 & 45 & 45 & 42 & 48 & 0.875 \\
4385 & 4386 & 44 & 48 & 44 & 48 & 0.916667 & 4384 & 4386 & 42 & 48 & 42 & 48 & 0.875 \\
3953 & 3954 & 34 & 35 & 33 & 36 & 0.916667 & 4383 & 4386 & 42 & 48 & 42 & 48 & 0.875 \\
21736 & 21739 & 66 & 70 & 65 & 71 & 0.915493 & 102964 & 102966 & 36 & 39 & 35 & 40 & 0.875 \\
49272 & 49273 & 34 & 33 & 32 & 35 & 0.914286 & 4380 & 4381 & 36 & 39 & 35 & 40 & 0.875 \\
3951 & 3953 & 33 & 34 & 32 & 35 & 0.914286 & 56480 & 56482 & 30 & 30 & 28 & 32 & 0.875 \\
3949 & 3950 & 32 & 33 & 31 & 34 & 0.911765 & 19322 & 19323 & 86 & 77 & 76 & 87 & 0.873563 \\
7277 & 7278 & 121 & 116 & 113 & 124 & 0.91129 & 21739 & 21740 & 70 & 78 & 69 & 79 & 0.873418 \\
4384 & 4385 & 42 & 44 & 41 & 45 & 0.911111 & 14473 & 14474 & 56 & 62 & 55 & 63 & 0.873016 \\
56485 & 56486 & 51 & 56 & 51 & 56 & 0.910714 & 49272 & 49274 & 34 & 39 & 34 & 39 & 0.871795 \\
40051 & 40052 & 72 & 75 & 70 & 77 & 0.909091 & 42969 & 42974 & 56 & 60 & 54 & 62 & 0.870968 \\
104123 & 104124 & 40 & 44 & 40 & 44 & 0.909091 & 14226 & 14227 & 51 & 50 & 47 & 54 & 0.87037 \\
15127 & 15128 & 62 & 62 & 59 & 65 & 0.907692 & 4386 & 4387 & 48 & 53 & 47 & 54 & 0.87037 \\
14118 & 14119 & 81 & 83 & 78 & 86 & 0.906977 & 15535 & 15537 & 231 & 242 & 220 & 253 & 0.869565 \\
44669 & 44670 & 40 & 42 & 39 & 43 & 0.906977 & 14030 & 14119 & 89 & 83 & 80 & 92 & 0.869565 \\
58919 & 58920 & 71 & 72 & 68 & 75 & 0.906667 & 32574 & 32575 & 55 & 59 & 53 & 61 & 0.868852 \\
21846 & 21847 & 50 & 51 & 48 & 53 & 0.90566 & 40054 & 40056 & 92 & 93 & 86 & 99 & 0.868687 \\
21736 & 21738 & 66 & 71 & 65 & 72 & 0.902778 & 3953 & 3956 & 34 & 37 & 33 & 38 & 0.868421 \\
3955 & 3956 & 39 & 37 & 36 & 40 & 0.9 & 15533 & 15534 & 187 & 196 & 178 & 205 & 0.868293 \\
3956 & 3957 & 37 & 39 & 36 & 40 & 0.9 & 42701 & 42702 & 262 & 261 & 243 & 280 & 0.867857 \\
19321 & 19326 & 90 & 94 & 87 & 97 & 0.896907 & 5682 & 5684 & 123 & 131 & 118 & 136 & 0.867647 \\
11089 & 79749 & 46 & 45 & 43 & 48 & 0.895833 & 40055 & 40056 & 90 & 93 & 85 & 98 & 0.867347 \\
11810 & 125309 & 45 & 46 & 43 & 48 & 0.895833 & 39309 & 39312 & 117 & 122 & 111 & 128 & 0.867188 \\
3958 & 3960 & 44 & 47 & 43 & 48 & 0.895833 & 50992 & 50993 & 42 & 42 & 39 & 45 & 0.866667 \\
21735 & 21736 & 61 & 66 & 60 & 67 & 0.895522 & 13975 & 13976 & 256 & 274 & 246 & 284 & 0.866197 \\
21732 & 21733 & 53 & 55 & 51 & 57 & 0.894737 & 14478 & 14479 & 88 & 93 & 84 & 97 & 0.865979 \\
88814 & 88815 & 37 & 35 & 34 & 38 & 0.894737 & 13977 & 13978 & 289 & 295 & 271 & 313 & 0.865815 \\
3954 & 3956 & 35 & 37 & 34 & 38 & 0.894737 & 89051 & 89053 & 47 & 50 & 45 & 52 & 0.865385 \\
2352 & 2353 & 65 & 58 & 58 & 65 & 0.892308 & 89049 & 89053 & 47 & 50 & 45 & 52 & 0.865385 \\
3952 & 3954 & 35 & 35 & 33 & 37 & 0.891892 & 88813 & 88815 & 34 & 35 & 32 & 37 & 0.864865 \\
69779 & 69784 & 125 & 119 & 115 & 129 & 0.891473 & 6166 & 6168 & 52 & 58 & 51 & 59 & 0.864407 \\
58909 & 58910 & 44 & 43 & 41 & 46 & 0.891304 & 27579 & 59022 & 259 & 263 & 242 & 280 & 0.864286 \\
\bottomrule
\end{tabular}
\end{table*}

% =====================================================================
% Section 2: Popularity Baseline
% =====================================================================
\section{Popularity Baseline}

\paragraph{Popularity Score and Recommendation Generation.} The popularity baseline recommends the same globally popular books to every user, personalized only by removing books the user has already seen. We defined popularity as the number of training interactions with \texttt{is\_read = 1} for each book, matching the implicit consumption signal that is the most abundant feedback type in the dataset. This choice also makes the baseline directly comparable to the implicit ALS model, since both use the same behavioral signal.

To generate recommendations, we first selected the top 1{,}000 globally popular books. Retaining a candidate pool larger than the final top-100 provides headroom: heavy readers may have already seen many popular books, and without extra candidates they could receive fewer than 100 recommendations. For each training user, we removed all books present in their training history via a left-anti join, then retained the top 100 unseen books ranked by global read count. Recommendations were computed once over all training users and cached, so the same precomputed list was reused across both label definitions and both evaluation splits.

\paragraph{Evaluation.} We evaluated under two label definitions to make the baseline comparable with later models. The first treats \texttt{is\_read = 1} as relevance, consistent with the popularity signal. The second treats \texttt{rating $\geq$ 4} as relevance, matching the preference signal used by the explicit ALS model.

Under \texttt{is\_read = 1} labels, the baseline achieved MAP 0.0320, NDCG@100 0.1081, precision@100 0.0275, and recall@100 0.1607 on the validation set, and MAP 0.0322, NDCG@100 0.1087, precision@100 0.0276, recall@100 0.1601 on the test set. Under \texttt{rating $\geq$ 4} labels, validation performance was MAP 0.0341, NDCG@100 0.1071, precision@100 0.0206, recall@100 0.1770, and test performance was MAP 0.0344, NDCG@100 0.1073, precision@100 0.0206, recall@100 0.1754. The close agreement between validation and test scores across both label definitions suggests the baseline is stable and not sensitive to the particular split.

\begin{table}[!ht]
\centering
\caption{Popularity baseline evaluation results on the 3\% sample.}
\label{tab:popularity}
\footnotesize
\setlength{\tabcolsep}{3pt}
\begin{tabular}{l l r r r r}
\toprule
\textbf{Split} & \textbf{Label} & \textbf{MAP} & \textbf{NDCG@100} & \textbf{P@100} & \textbf{R@100} \\
\midrule
Val.  & is\_read=1     & 0.032 & 0.1081 & 0.0275 & 0.1607 \\
Test  & is\_read=1     & 0.0322 & 0.1087 & 0.0276 & 0.1601 \\
Val.  & rating $\geq$ 4 & 0.0341 & 0.1071 & 0.0206 & 0.177 \\
Test  & rating $\geq$ 4 & 0.0344 & 0.1073 & 0.0206 & 0.1754 \\
\bottomrule
\end{tabular}
\end{table}

% =====================================================================
% Section 3: Explicit-Feedback ALS Recommender
% =====================================================================
\section{Explicit-Feedback ALS Recommender}

\paragraph{Model and evaluation setup.} We trained Spark ALS using only the numerical Goodreads \texttt{rating} field. Rows with \texttt{rating = 0} were excluded because they indicate missing explicit feedback rather than negative preference; \texttt{is\_read} and \texttt{is\_reviewed} were not used as inputs. We used \texttt{coldStartStrategy="drop"} and evaluated top-100 recommendations with Spark MLlib's \texttt{RankingMetrics}, after removing books already present in each user's training history. Relevance was defined as held-out \texttt{rating $\geq$ 4}, matching the explicit preference signal and the rating-label popularity baseline.

\paragraph{Data and tuning process.} We first screened ALS hyperparameters on a deterministic 1\% user sample to keep the grid affordable, then carried the selected configuration to the deterministic 3\% user sample used for the final explicit-ALS comparison. The 3\% sample contained 6{,}831{,}620 interactions from 26{,}335 users before filtering. After removing users with fewer than five interactions, 6{,}827{,}622 interactions and 24{,}415 users remained, split into 4{,}086{,}857 train rows, 1{,}365{,}418 validation rows, and 1{,}375{,}347 test rows.

\paragraph{Hyperparameter screen.} The 1\% validation screen covered \texttt{rank} $\in \{5, 20, 50, 100\}$ and \texttt{regParam} $\in \{0.001, 0.01, 0.1, 0.5\}$, with \texttt{recPool=250} and validation NDCG@100 as the selection criterion. We evaluated 13 of 16 cells; the skipped cells were unlikely to change selection after neighboring low-rank and heavily regularized settings performed poorly. Table~\ref{tab:sweep} shows the evaluated cells, sorted by NDCG@100.

\begin{table}[!ht]
\centering
\caption{1\% validation grid for explicit ALS (\texttt{recPool=250}, \texttt{maxIter=5}). Sorted by NDCG@100.}
\label{tab:sweep}
\footnotesize
\setlength{\tabcolsep}{3pt}
\begin{tabular}{r r r r r r}
\toprule
\textbf{Rank} & \textbf{regParam} & \textbf{NDCG@100} & \textbf{MAP} & \textbf{P@100} & \textbf{R@100} \\
\midrule
100 & 0.01  & \textbf{0.014463} & 0.00281 & 0.003379 & 0.034114 \\
100 & 0.001 & 0.014431          & 0.002534 & 0.003807 & 0.032669 \\
50  & 0.01  & 0.005865          & 0.001026 & 0.00174 & 0.013486 \\
50  & 0.001 & 0.005863          & 0.000665 & 0.001933 & 0.013017 \\
50  & 0.1   & 0.002486          & 0.000415 & 0.000567 & 0.006412 \\
100 & 0.1   & 0.002212          & 0.000401 & 0.000445 & 0.005932 \\
20  & 0.01  & 0.002061          & 0.000232 & 0.00077 & 0.004884 \\
20  & 0.1   & 0.001412          & 0.000171 & 0.000458 & 0.003597 \\
20  & 0.5   & 0.000834          & 0.000107 & 0.000197 & 0.002278 \\
50  & 0.5   & 0.000809          & 0.0001 & 0.000209 & 0.002037 \\
5   & 0.01  & 0.000172          & 0.00001 & 0.000086 & 0.000345 \\
5   & 0.1   & 0.000168          & 0.00001 & 0.000074 & 0.000438 \\
5   & 0.5   & 0.000154          & 0.000015 & 0.000042 & 0.000496 \\
\bottomrule
\end{tabular}
\end{table}

Increasing \texttt{rank} from 5 to 100 raises NDCG@100 by roughly two orders of magnitude, indicating that the model lacks sufficient capacity at low rank and benefits substantially from more latent dimensions on this sample. The top two configurations were \texttt{rank=100}, \texttt{regParam=0.01} (NDCG@100 $=0.014463$) and \texttt{rank=100}, \texttt{regParam=0.001} (NDCG@100 $=0.014431$). These scores differ by only $0.000032$, so they are effectively tied; we selected \texttt{regParam=0.01} because it had the highest validation score and provided slightly stronger regularization. Performance collapses at \texttt{regParam=0.1}, suggesting that overly strong regularization erases the latent structure that high rank otherwise recovers.

\paragraph{3\% validation and test.} Since the 1\% grid selected \texttt{rank=100}, \texttt{regParam=0.01}, the 3\% experiment reran this setting with candidate recommendation pools \texttt{recPool} $\in \{250, 500\}$ and included lower-rank sanity checks, all with \texttt{maxIter=5}. The selected high-rank setting remained best on validation.

\begin{table}[!ht]
\centering
\caption{3\% validation comparison for the selected ALS hyperparameters (\texttt{rank=100}, \texttt{regParam=0.01}).}
\label{tab:sweep3pct}
\footnotesize
\setlength{\tabcolsep}{3pt}
\begin{tabular}{r r r r r}
\toprule
\textbf{recPool} & \textbf{MAP} & \textbf{NDCG@100} & \textbf{P@100} & \textbf{R@100} \\
\midrule
250 & 0.001363 & 0.006217          & 0.001431 & 0.014823 \\
500 & 0.001368 & \textbf{0.006303} & 0.001529 & 0.01489 \\
\bottomrule
\end{tabular}
\end{table}

Using validation NDCG@100 as the criterion, \texttt{recPool=500} was slightly better than \texttt{recPool=250}; its validation RMSE was 1.683. On the 3\% test split, the selected setting (\texttt{rank=100}, \texttt{regParam=0.01}, \texttt{recPool=500}, \texttt{maxIter=5}) achieved MAP $=$ 0.00158, NDCG@100 $=$ 0.00678, P@100 $=$ 0.00158, and R@100 $=$ 0.01583. RMSE was 1.679, but ranking metrics were the basis for model selection and comparison.

\begin{table}[!ht]
\centering
\caption{3\% test-set comparison using the shared \texttt{rating $\geq$ 4} relevance label.}
\label{tab:compare3pct}
\footnotesize
\setlength{\tabcolsep}{3pt}
\begin{tabular}{l r r r r}
\toprule
\textbf{Model} & \textbf{MAP} & \textbf{NDCG@100} & \textbf{P@100} & \textbf{R@100} \\
\midrule
Popularity & 0.034398 & 0.107278 & 0.020638 & 0.175361 \\
Expl.~ALS  & 0.001581 & 0.006778 & 0.001584 & 0.015833 \\
\bottomrule
\end{tabular}
\end{table}

The explicit ALS model underperformed the 3\% popularity baseline by a wide margin under the same \texttt{rating $\geq$ 4} label. This suggests that the explicit rating graph remained too sparse at the sampled scale for ALS to beat the heavy-tailed global popularity signal. The small gain from \texttt{recPool=250} to 500 also suggests that candidate-pool truncation was not the main bottleneck. We therefore interpret this result as an empirical finding about the sampled explicit-rating signal, not as evidence that the ALS implementation failed.

% =====================================================================
% Section 4: Implicit-Feedback ALS Recommender
% =====================================================================
\section{Implicit-Feedback ALS Recommender}

\paragraph{Implicit Feedback Signal.} For this model we constructed a composite behavioral signal, deliberately excluding numerical ratings. We defined an implicit score as
\[
\text{implicit\_score} = \text{is\_read} + 2 \times \text{is\_reviewed}.
\]
The coefficient of 2 on \texttt{is\_reviewed} reflects the intuition that writing a review demands more active engagement than simply marking a book as read, so reviewed books are treated as stronger evidence of interest. Interactions where both fields are zero were dropped, as they carry no behavioral signal. Under this scheme, a book that was read but not reviewed receives a score of 1; a book that is both read and reviewed receives a score of 3. Any zero-score row is excluded from training entirely.

\paragraph{Model Training and Hyperparameter Tuning.} We used Spark ALS with \texttt{implicitPrefs=True}. In this mode, Spark follows the standard implicit-feedback ALS approach, where positive interactions are treated as preferences and larger interaction values receive higher confidence. As with explicit ALS, we first screened rank and regularization on the deterministic 1\% user sample, then carried the selected configuration to the 3\% sample for final validation and test evaluation. The screen covered ranks $\{5, 20, 50, 100\}$ and regularization values $\{0.01, 0.1, 0.5\}$ while holding \texttt{alpha = 15.0} and \texttt{maxIter = 5} fixed. For evaluation, we removed books already seen in training via a left-anti join and scored the remaining top 100 against held-out \texttt{is\_read = 1} labels.

\begin{table}[!ht]
\centering
\caption{1\% validation grid for implicit ALS (\texttt{alpha = 15.0}, \texttt{maxIter = 5}). Sorted by NDCG@100.}
\label{tab:implicit-grid}
\footnotesize
\setlength{\tabcolsep}{3pt}
\begin{tabular}{r r r r r r}
\toprule
\textbf{Rank} & \textbf{regParam} & \textbf{NDCG@100} & \textbf{MAP} & \textbf{P@100} & \textbf{R@100} \\
\midrule
100 & 0.1  & \textbf{0.171637} & 0.056755 & 0.045476 & 0.250553 \\
100 & 0.01 & 0.170727 & 0.056483 & 0.044781 & 0.249645 \\
100 & 0.5  & 0.170672 & 0.056003 & 0.045988 & 0.248146 \\
50  & 0.5  & 0.151277 & 0.048491 & 0.039392 & 0.22201 \\
50  & 0.1  & 0.150733 & 0.048228 & 0.038992 & 0.222169 \\
50  & 0.01 & 0.149377 & 0.047633 & 0.038376 & 0.22143 \\
20  & 0.5  & 0.144677 & 0.045052 & 0.038533 & 0.21183 \\
20  & 0.1  & 0.143191 & 0.044465 & 0.037988 & 0.210762 \\
20  & 0.01 & 0.139779 & 0.043105 & 0.037109 & 0.206289 \\
5   & 0.1  & 0.119108 & 0.035218 & 0.03206 & 0.174895 \\
5   & 0.5  & 0.1186 & 0.035038 & 0.031962 & 0.17405 \\
5   & 0.01 & 0.117414 & 0.034639 & 0.031375 & 0.173287 \\
\bottomrule
\end{tabular}
\end{table}

Rank 100 with \texttt{regParam = 0.1} was the best configuration, achieving NDCG@100 of about 0.1716 on validation. The top rank-100 settings were tightly clustered: \texttt{regParam=0.1} beat \texttt{regParam=0.01} by about 0.0009 NDCG@100, and \texttt{regParam=0.5} was only slightly lower. We therefore treat the choice of \texttt{regParam=0.1} as a small validation-set advantage rather than a practically large effect; repeated user-history splits could plausibly change this ordering. The effect of regularization depended on rank: at ranks 20 and 50, higher regularization (0.5) performed best, while at rank 100 a moderate value (0.1) was optimal. The lowest rank (5) produced the weakest results regardless of regularization, consistent with the expectation that a large, sparse dataset benefits from richer latent representations.

\paragraph{3\% validation and test.} We then reran the selected configuration on the 3\% user sample with \texttt{recPool=500}. The selected configuration was stable but not exactly identical across sample and candidate-pool settings: validation NDCG@100 was 0.171637 on the 1\% sample with \texttt{recPool=250}, 0.171635 on the 3\% sample with \texttt{recPool=250}, and 0.171636 on the 3\% sample with \texttt{recPool=500}. Thus the apparent equality at four decimal places is a rounding effect, not literally identical outputs. For implicit ALS, \texttt{recPool=250} was already enough to generate stable top-100 recommendations after filtering out seen books, so increasing \texttt{recPool} to 500 changed the metrics only in later decimals; explicit ALS can behave differently because it uses much sparser rating-only data, so expanding the candidate pool can slightly change the final top-100 list. On the 3\% test set with \texttt{recPool=500}, results were MAP 0.0566, NDCG@100 0.1722, precision@100 0.0458, recall@100 0.2509 --- nearly identical to validation, indicating the selected hyperparameters generalize without overfitting to the validation split.

\begin{table}[!ht]
\centering
\caption{Implicit ALS final 3\% evaluation (\texttt{rank = 100}, \texttt{regParam = 0.1}, \texttt{alpha = 15.0}, \texttt{recPool=500}).}
\label{tab:implicit-final}
\footnotesize
\setlength{\tabcolsep}{3pt}
\begin{tabular}{l r r r r}
\toprule
\textbf{Split} & \textbf{MAP} & \textbf{NDCG@100} & \textbf{P@100} & \textbf{R@100} \\
\midrule
Val.  & 0.056755 & 0.171636 & 0.045476 & 0.250547 \\
Test  & 0.05664 & 0.172238 & 0.045846 & 0.250926 \\
\bottomrule
\end{tabular}
\end{table}

\paragraph{Comparison with Popularity Baseline.} Compared with the popularity baseline on the test set under the same \texttt{is\_read = 1} labels (MAP 0.0322, NDCG@100 0.1087, precision@100 0.0276, recall@100 0.1601), the implicit ALS model yields substantial gains: approximately 76\% relative improvement in MAP, 58\% in NDCG@100, 66\% in precision@100, and 57\% in recall@100. These results confirm that personalization --- even without any explicit preference signal --- substantially outperforms a non-personalized global ranking.

\paragraph{Comparison with Explicit ALS.} Compared with the explicit-feedback ALS model from Section~3, implicit ALS achieved much stronger ranking metrics, though under a different relevance label. Explicit ALS on the 3\% test split reached MAP 0.001581 and NDCG@100 0.006778 using \texttt{rating $\geq$ 4} labels, while implicit ALS reached MAP 0.0566 and NDCG@100 0.1722 using \texttt{is\_read = 1} labels. This gap should not be interpreted as a perfectly controlled model comparison because the evaluation labels differ; rather, it shows that implicit ALS is much better aligned with the abundant read-behavior signal than explicit ALS is with the sparse rating graph.

\paragraph{Limitations and potential fixes.} The implicit model treats reading and reviewing as positive behavioral evidence, but neither signal guarantees that the user liked the book. This can inflate performance under an \texttt{is\_read = 1} label while still recommending books that were consumed but not preferred. The review weight was also fixed at 2 rather than tuned jointly with \texttt{alpha}, so future runs should search over confidence weights, \texttt{alpha}, and \texttt{maxIter}. Finally, evaluating implicit ALS under both read-based and rating-based relevance labels would make its tradeoff against explicit ALS clearer.

% =====================================================================
% Section 5: Combined Explicit + Implicit ALS
% =====================================================================
\section{Combined Explicit and Implicit ALS}

\paragraph{Score transformation.} The combined model trains on a single \texttt{combined\_score} column built from rating, read, and review signals. For user $u$ and book $i$,
\[
\text{combined\_score}_{ui} =
\begin{cases}
r_{ui} & r_{ui} > 0 \\
2.0    & r_{ui} = 0,\ v_{ui} = 1 \\
1.0    & r_{ui} = 0,\ v_{ui} = 0,\ d_{ui} = 1 \\
\text{excluded} & \text{otherwise.}
\end{cases}
\]
where $d_{ui}$ is \texttt{is\_read} and $v_{ui}$ is \texttt{is\_reviewed}. Ratings are kept unchanged; read/review signals are used only when no rating is available. This avoids treating a low-rated book as a strong positive simply because it was read or reviewed. This follows the same general intuition as confidence weighting in implicit-feedback recommenders: higher-effort actions provide stronger evidence than passive consumption. ALS was run in explicit mode (\texttt{implicitPrefs = False}) with \texttt{ratingCol = combined\_score} and top-100 evaluation.

\paragraph{Sampling and splitting.} The final experiment uses the deterministic 3\% user sample and the shared per-user 60/20/20 split described in the intro (seed 42, users with fewer than 5 interactions removed). Keeping validation and test users connected to training data is necessary for ALS recommendations.

\paragraph{Cluster-conscious design.} The implementation is staged to reduce reruns on Dataproc: the 3\% sample and split are saved once, each trained ALS model is persisted to HDFS, validation and test metrics are saved as CSV, and saved models are reused when only evaluation needs to be rerun. The grid was kept intentionally small rather than running a wide hyperparameter sweep, in line with the assignment guidance to sample when needed, save intermediate outputs, and avoid unnecessary recomputation.

\paragraph{Process justification.} We first used the completed 1\% runs to choose the combined-score weights, then carried the best weights to the 3\% sample and tuned only \texttt{rank} and \texttt{regParam}. This keeps the search goal-oriented: pick a reasonable signal transformation cheaply, then spend the larger 3\% budget on the ALS parameters most likely to affect ranking quality.

\paragraph{Weight selection (1\% sample).} The 1\% weight grid was run at \texttt{rank=50}, \texttt{regParam=0.1}, \texttt{recPool=250}. Table~\ref{tab:combined-weights} shows that the lightest weights $(1.0, 2.0)$ won on NDCG@100, so these were carried to the 3\% sweep.

\begin{table}[!ht]
\centering
\caption{1\% weight grid for combined ALS (\texttt{rank=50}, \texttt{regParam=0.1}, \texttt{recPool=250}). Best NDCG@100 in bold.}
\label{tab:combined-weights}
\footnotesize
\setlength{\tabcolsep}{3pt}
\begin{tabular}{r r r r r r}
\toprule
\textbf{Read} & \textbf{Review} & \textbf{MAP} & \textbf{NDCG@100} & \textbf{P@100} & \textbf{R@100} \\
\midrule
1 & 2 & 0.000538 & \textbf{0.003113} & 0.000702 & 0.008151 \\
2 & 3 & 0.000526 & 0.0026          & 0.000548 & 0.006694 \\
3 & 4 & 0.000702 & 0.002471          & 0.000451 & 0.005923 \\
\bottomrule
\end{tabular}
\end{table}

\paragraph{Rank/regularization sweep (3\% sample).} With the selected weights, we tuned \texttt{rank} and \texttt{regParam} on the 3\% sample at \texttt{recPool=500}, \texttt{maxIter=5} (Table~\ref{tab:combined-grid}).

\begin{table}[!ht]
\centering
\caption{3\% rank/regularization grid for combined ALS (\texttt{read\_weight=1.0}, \texttt{review\_weight=2.0}, \texttt{recPool=500}, \texttt{maxIter=5}). Best NDCG@100 in bold.}
\label{tab:combined-grid}
\footnotesize
\setlength{\tabcolsep}{3pt}
\begin{tabular}{r r r r r r}
\toprule
\textbf{Rank} & \textbf{regParam} & \textbf{MAP} & \textbf{NDCG@100} & \textbf{P@100} & \textbf{R@100} \\
\midrule
100 & 0.01 & 0.001654 & \textbf{0.007444} & 0.001874 & 0.01705 \\
100 & 0.1  & 0.000368 & 0.001925          & 0.000384 & 0.005151 \\
50  & 0.1  & 0.000311 & 0.001623          & 0.000357 & 0.004057 \\
\bottomrule
\end{tabular}
\end{table}

\paragraph{Selected configuration.} \texttt{rank=100}, \texttt{regParam=0.01}, \texttt{read\_weight=1.0}, \texttt{review\_weight=2.0}, \texttt{recPool=500}, \texttt{maxIter=5}.

\paragraph{Final results.} The primary relevance label is \texttt{rating $\geq$ 4}; we additionally report the test split under \texttt{is\_read = 1} for completeness (Table~\ref{tab:combined-final}).

\begin{table}[!ht]
\centering
\caption{Combined ALS final evaluation on the 3\% sample.}
\label{tab:combined-final}
\footnotesize
\setlength{\tabcolsep}{3pt}
\begin{tabular}{l l r r r r}
\toprule
\textbf{Split} & \textbf{Label} & \textbf{MAP} & \textbf{NDCG@100} & \textbf{P@100} & \textbf{R@100} \\
\midrule
Val.  & rating $\geq$ 4 & 0.001654 & 0.007444 & 0.001874 & 0.01705 \\
Test  & rating $\geq$ 4 & 0.001751 & 0.007562 & 0.001859 & 0.01704 \\
Test  & is\_read = 1    & 0.001616 & 0.007505 & 0.002169 & 0.015577 \\
\bottomrule
\end{tabular}
\end{table}

\paragraph{Comparison with baseline and other models.} Table~\ref{tab:combined-compare} compares the combined model against the popularity baseline and explicit ALS model on the 3\% test split under the shared \texttt{rating $\geq$ 4} relevance label. Table~\ref{tab:combined-read-compare} separately compares read-label performance against implicit ALS, because implicit ALS was evaluated under \texttt{is\_read = 1}.

\begin{table}[!ht]
\centering
\caption{3\% test-set comparison using the shared \texttt{rating $\geq$ 4} relevance label.}
\label{tab:combined-compare}
\footnotesize
\setlength{\tabcolsep}{3pt}
\begin{tabular}{l r r r r}
\toprule
\textbf{Model} & \textbf{MAP} & \textbf{NDCG@100} & \textbf{P@100} & \textbf{R@100} \\
\midrule
Popularity      & 0.034398 & 0.107278 & 0.020638 & 0.175361 \\
Expl.~ALS       & 0.001581 & 0.006778 & 0.001584 & 0.015833 \\
Combined ALS    & 0.001751 & 0.007562 & 0.001859 & 0.01704 \\
\bottomrule
\end{tabular}
\end{table}

\begin{table}[!ht]
\centering
\caption{3\% test-set comparison using the shared \texttt{is\_read = 1} relevance label.}
\label{tab:combined-read-compare}
\footnotesize
\setlength{\tabcolsep}{3pt}
\begin{tabular}{l r r r r}
\toprule
\textbf{Model} & \textbf{MAP} & \textbf{NDCG@100} & \textbf{P@100} & \textbf{R@100} \\
\midrule
Popularity   & 0.0322 & 0.1087 & 0.0276 & 0.1601 \\
Implicit ALS & 0.05664 & 0.172238 & 0.045846 & 0.250926 \\
Combined ALS & 0.001616 & 0.007505 & 0.002169 & 0.015577 \\
\bottomrule
\end{tabular}
\end{table}

Among the rating-based models, the combined ALS edges out explicit-only ALS on every metric, suggesting that adding read/review signals as fallback evidence when ratings are missing recovers some additional ranking signal. Both, however, underperform the popularity baseline by roughly an order of magnitude. Under the read-label evaluation, implicit ALS is much stronger than combined ALS because its training objective and evaluation label are directly aligned, while the combined model remains an explicit-mode ALS model trained on a rating-like score.

\paragraph{Limitations and potential fixes.} The combined model uses 3\% of users; full-data training could improve overlap but would require more cluster resources. \texttt{recPool=500} may still be too small after filtering seen books; raising it to 1000 could help recall if the cluster can sustain the cost. The grid was deliberately sparse, and exploring more rank/regularization/iteration values would likely improve results. The score transformation itself is simple; an alternative is to treat \texttt{rating} as the explicit preference signal while using \texttt{is\_read} and \texttt{is\_reviewed} as per-interaction confidence modifiers (Hu--Koren--Volinsky style). Finally, a hybrid model that blends ALS with a popularity prior may work better given how strong the popularity baseline is on this dataset.

% =====================================================================
% Contributions and software
% =====================================================================
\section*{Contributions and Software}

\paragraph{Member contributions.} All three members contributed jointly to data preprocessing, the train/validation/test split, the five deliverables, and the writing of this report. Yoonseol Jang focused on the explicit-feedback ALS experiments, Donggyu Kim focused on the implicit-feedback ALS experiments, and David Lee focused on the combined-feedback experiments; all members participated across each stage and reviewed the final results.

\paragraph{AI contributions.} Generative AI was used for grammar checking, table creation, text formatting, and code refactoring. All experimental design decisions, hyperparameter choices, and result interpretations were made by the team.

\paragraph{Additional software.} All experiments used Apache Spark (Spark SQL, Spark MLlib's \texttt{ALS}, \texttt{MinHashLSH}, and \texttt{RankingMetrics}). No other libraries beyond the standard Spark distribution were required.

\end{document}
